% =====================================================================
%  YUYOS DE LA QUEBRADA DE CAFAYATE — CUBIERTA PARA IMPRENTA
%
%  Un solo pliego: contratapa | lomo | tapa, con 3 mm de demasía en
%  todo el contorno. Compilar con:  pdflatex cubierta.tex
%
%  Antes de mandar a imprenta, confirmar con la imprenta los dos datos
%  de abajo (páginas del interior y espesor de la hoja del papel
%  elegido). El lomo se calcula solo.
% =====================================================================

% ------------------------------ Datos --------------------------------
\def\paginas{60}          % páginas del interior (pdfinfo cafayate.pdf)
\def\espesorhoja{0.11}    % mm por HOJA (dos páginas); pedirlo a la imprenta
\def\demasia{3}           % mm
\def\anchotapa{170}       % mm
\def\altotapa{240}        % mm

\documentclass{article}
\usepackage[utf8]{inputenc}
\usepackage[T1]{fontenc}
\usepackage[spanish]{babel}
\usepackage{mathpazo}
\renewcommand{\sfdefault}{pag}
\usepackage[cmyk]{xcolor}
% Láminas en CMYK y perfil de salida (OutputIntent), generados por cmyk.py.
\IfFileExists{perfil-salida.tex}{%
  \newcommand{\carpetaimg}{img-cmyk}%
  % perfil-salida.tex — generado por cmyk.py. No editar a mano.
% Declara el perfil de salida (OutputIntent) del PDF de imprenta.
% Perfil: FOGRA39L Coated
\immediate\pdfobj stream attr{/N 4} file{perfil-salida.icc}
\pdfcatalog{/OutputIntents [ <<
  /Type /OutputIntent
  /S /GTS_PDFX
  /OutputCondition (FOGRA39L Coated)
  /OutputConditionIdentifier (FOGRA39)
  /RegistryName (http://www.color.org)
  /Info (FOGRA39L Coated)
  /DestOutputProfile \the\pdflastobj\space 0 R
>> ]}
%
}{%
  \newcommand{\carpetaimg}{img}%
  \PackageWarningNoLine{cubierta}{Falta perfil-salida.tex: correr
    "py cmyk.py" antes de compilar}%
}
\usepackage{graphicx}
\usepackage{tikz}
\usetikzlibrary{fadings}

% Lomo en mm, redondeado a décimas
\pgfmathsetmacro{\lomo}{round(\paginas/2*\espesorhoja*10)/10}
\pgfmathsetmacro{\anchototal}{2*\anchotapa+\lomo+2*\demasia}
\pgfmathsetmacro{\altototal}{\altotapa+2*\demasia}

\usepackage[paperwidth=\anchototal mm, paperheight=\altototal mm, margin=0pt]{geometry}
% Caja de corte (TrimBox) y de demasía (BleedBox), en puntos PostScript
\pgfmathsetmacro{\dbp}{\demasia*72/25.4}
\pgfmathsetmacro{\wbp}{\anchototal/25.4*72}
\pgfmathsetmacro{\hbp}{\altototal/25.4*72}
\pgfmathsetmacro{\tw}{\wbp-\dbp}
\pgfmathsetmacro{\th}{\hbp-\dbp}
% (se aplica al comenzar el documento)

\definecolor{verdehuerta}{HTML}{3F4A2E}
\definecolor{ocre}       {HTML}{C28A34}
\definecolor{crema}      {HTML}{FBF4E6}
\pagestyle{empty}

\begin{document}
\edef\cajas{/TrimBox [\dbp\space\dbp\space\tw\space\th] /BleedBox [0 0 \wbp\space\hbp]}%
\pdfpageattr\expandafter{\cajas}
\noindent
\begin{tikzpicture}[x=1mm, y=1mm, remember picture, overlay]
  % Coordenadas: (0,0) es el borde inferior izquierdo con demasía.
  \pgfmathsetmacro{\xlomo}{\demasia+\anchotapa}           % empieza el lomo
  \pgfmathsetmacro{\xtapa}{\xlomo+\lomo}                   % empieza la tapa
  \pgfmathsetmacro{\anchopanel}{\anchotapa+\demasia}
  \coordinate (O) at (current page.south west);

  % ---------------- Contratapa (izquierda) ----------------
  \begin{scope}[shift={(O)}]
    \clip (0,0) rectangle (\anchopanel,\altototal);
    \node[anchor=center, inner sep=0] at (\anchopanel/2,\altototal/2)
      {\includegraphics[width=\anchopanel mm]{\carpetaimg/23-contratapa.jpg}};
    \fill[verdehuerta, opacity=0.72] (0,\altototal) rectangle (\anchopanel,\altototal-\demasia-120);
    \fill[verdehuerta, opacity=0.72, path fading=south]
      (0,\altototal-\demasia-120) rectangle (\anchopanel,\altototal-\demasia-165);
    \node[anchor=north, inner sep=0] at (\demasia+\anchotapa/2,\altototal-\demasia-28)
      {\begin{minipage}{130mm}\centering
        {\sffamily\fontsize{13}{17}\selectfont\color{crema} YUYOS DE LA QUEBRADA}\\[0.6em]
        {\color{ocre}\rule{3.4cm}{0.8pt}}\\[0.9em]
        {\itshape\small\color{crema}
         «Existe una idea antigua que dice que en un perímetro de un
         kilómetro a la redonda tenemos todo lo que necesitamos para
         sanarnos. Por eso considero al ser urbano huérfano de madre
         naturaleza.»}\\[1.0em]
        {\sffamily\fontsize{9.5}{12}\selectfont\color{crema} MARÍA EUGENIA ORELLANA}\\[1.2em]
        {\small\color{crema}% Texto de contratapa. Lo usan cafayate.tex (modo pantalla) y cubierta.tex.
Chañar para la tos, jarilla para el reuma, cedrón para el ánimo. En la
Quebrada de Cafayate la botica está a un kilómetro a la redonda, y este
libro la recorre planta por planta: cómo reconocerla, qué parte se usa,
cómo se prepara y qué dice de ella la tradición. Veintiuna fichas nacidas
del cuadernillo de María Eugenia Orellana y de las charlas con los
vecinos, láminas para identificar cada especie en el campo y
advertencias claras sobre las plantas que no se tocan.%
}\\[1.0em]
        {\sffamily\fontsize{8}{11}\selectfont\color{crema} VALLES CALCHAQUÍES · SALTA}
      \end{minipage}};
    % Lugar del código de barras del ISBN (blanco, 50 x 30 mm)
    \fill[white] (\demasia+12,\demasia+12) rectangle (\demasia+62,\demasia+42);
    \node[font=\sffamily\tiny, text=black] at (\demasia+37,\demasia+27) {ISBN / código de barras};
  \end{scope}

  % ---------------- Lomo ----------------
  \begin{scope}[shift={(O)}]
    \fill[verdehuerta] (\xlomo,0) rectangle (\xtapa,\altototal);
    % Con menos de 6 mm de lomo el texto no entra legible: se omite.
    \pgfmathparse{\lomo>=6 ? 1 : 0}
    \ifnum\pgfmathresult=1
    \node[rotate=-90, font=\sffamily\fontsize{8}{10}\selectfont, text=crema]
      at ({(\xlomo+\xtapa)/2},\altototal/2)
      {YUYOS DE LA QUEBRADA DE CAFAYATE \quad{\color{ocre}·}\quad MARÍA EUGENIA ORELLANA};
    \fi
  \end{scope}

  % ---------------- Tapa (derecha) ----------------
  \begin{scope}[shift={(O)}]
    \clip (\xtapa,0) rectangle (\anchototal,\altototal);
    \node[anchor=center, inner sep=0] at ({\xtapa+\anchopanel/2},\altototal/2)
      {\includegraphics[width=\anchopanel mm]{\carpetaimg/01-cubierta.jpg}};
    \fill[verdehuerta, opacity=0.74, path fading=south]
      (\xtapa,\altototal) rectangle (\anchototal,\altototal-\demasia-106);
    \fill[verdehuerta, opacity=0.88]
      (\xtapa,0) rectangle (\anchototal,\demasia+34);
    \fill[verdehuerta, opacity=0.88, path fading=north]
      (\xtapa,\demasia+34) rectangle (\anchototal,\demasia+60);
    \node[anchor=north, inner sep=0] at ({\xtapa+\anchotapa/2},\altototal-\demasia-20)
      {\begin{minipage}{140mm}\centering
        {\sffamily\bfseries\fontsize{27}{31}\selectfont\color{crema} Yuyos}\\[0.25em]
        {\sffamily\fontsize{15}{19}\selectfont\color{crema} de la Quebrada de Cafayate}\\[0.6em]
        {\color{ocre}\rule{5cm}{1pt}}\\[0.7em]
        {\sffamily\fontsize{11}{14}\selectfont\color{crema} MARÍA EUGENIA ORELLANA}
      \end{minipage}};
    \node[anchor=south, inner sep=0] at ({\xtapa+\anchotapa/2},\demasia+15)
      {\begin{minipage}{140mm}\centering
        {\sffamily\bfseries\fontsize{9}{12}\selectfont\color{crema}
         LA BOTICA COMPLETA, A UN KILÓMETRO A LA REDONDA\\[0.3em]
         VALLES CALCHAQUÍES · SALTA}
      \end{minipage}};
  \end{scope}
\end{tikzpicture}
\end{document}
